\documentclass[12pt]{article}
\usepackage[margin=0.6in,top=1.15in,bottom=0.55in]{geometry}
\usepackage{lmodern}
\usepackage{microtype}
\usepackage{titlesec}
\usepackage{enumitem}
\usepackage[hidelinks]{hyperref}
\usepackage{../../latex/grader-worksheet}

\setlength{\parindent}{0pt}
\setlength{\parskip}{0.25em}
\titleformat{\section}{\large\bfseries\scshape}{}{0pt}{}

\GraderSetup{assignment-id=U1L8LE,template-revision=1}

\begin{document}

\begin{center}
{\LARGE\bfseries Week 8 Lit Example Worksheet}\par\vspace{0.05em}
{\large\scshape Macbeth: Fair, Foul, and Shifting Security}\par\vspace{0.12em}
\rule{0.78\textwidth}{0.5pt}
\end{center}

\textbf{Purpose:} Apply Week 8's method to \textit{Macbeth}. Expose shifting terms and retest the inference with fixed meanings.

\textbf{Directions:} Use the Lit Example Reader. Do not retell the whole plot.

\section*{Part 1: Fair and the Greetings}

\textbf{Question 1.1}\\[-0.15em]
In Passage 1, give two ordinary senses of \textit{fair} that the language can carry.

\GraderAnswerBox{Part 1 Q1}{2.15in}

\textbf{Question 1.2}\\[-0.15em]
Which greetings are present facts, and which are future claims? Why does that difference matter for a stable argument?

\GraderAnswerBox{Part 1 Q2}{2.35in}

\newpage

\section*{Part 2: From Sounding Fair to Acting}

\textbf{Question 2.1}\\[-0.15em]
Reconstruct the risky inference Macbeth may draw from the witches' words in a middle-term shape (even if informal). Name the term that is doing the linking work.

\GraderAnswerBox{Part 2 Q1}{2.45in}

\textbf{Question 2.2}\\[-0.15em]
Explain why that link is an ambiguous middle or term shift rather than a stable proof that he must be king by crime.

\GraderAnswerBox{Part 2 Q2}{2.35in}

\newpage

\section*{Part 3: Security Promises}

\textbf{Question 3.1}\\[-0.15em]
For ``none of woman born shall harm Macbeth,'' state Sense A (how Macbeth hears it) and Sense B (a different sense that later undercuts the guarantee).

\GraderAnswerBox{Part 3 Q1}{2.45in}

\textbf{Question 3.2}\\[-0.15em]
For the Birnam wood promise, state Sense A and Sense B, then say whether Macbeth's ``That will never be'' still follows under both fixed senses.

\GraderAnswerBox{Part 3 Q2}{2.45in}

\newpage

\section*{Part 4: Repair}

\textbf{Question 4.1}\\[-0.15em]
Rewrite one of the security arguments with fixed meanings so that it no longer falsely guarantees Macbeth's safety.

\GraderAnswerBox{Part 4 Q1}{2.55in}

\textbf{Question 4.2}\\[-0.15em]
In one paragraph, explain how Week 8's tool changes a reading of these passages from ``magic prophecy'' to ``unstable terms used as if they were proof.''

\GraderAnswerBox{Part 4 Q2}{2.45in}

\end{document}
